\section{讲解笔记：【第1次作业笔记】 几何组成分析(1)}
\begin{analysisbox}
本节按题目、答案和讲解笔记顺序整理。对原图中的结构几何关系，正文保留 可辨识文字，并在图形页中保留清理后的题图，便于核对杆件、支座、荷载和尺寸。
\end{analysisbox}
\subsection*{第 1 页转写}
结力基础班第1次作业一一几何组成分析

几何组成分析作业及答案

1-1图1-56所示体系的几何组成为（

）。（浙江大学2008）

A.几何瞬变，无多余约束

B.几何不变

C.几何常变

D.几何瞬变，有多余约束

瞬变体系，两个多余约束。

平行且等长，是常变。

平行不等长，是瞬变。

图1-56

1-2欲使图1-57所示体系成为无多余约束的几何不变体系，则需在A端加人（B）。

（中南大学2011，合肥工业大学2009）

B.固定支座

A.固定铰支座

C.滑动铰支座

D.定向支座

二元体

图1-57

（1-3图1-58所示体系的几何构造性质是（C）。（中南大学2009)

A.常变体系

B.瞬变体系

C.无多余联系的几何不变体系

D.有多余联系的几何不变体系

图1-58

\subsection*{第 2 页转写}
1-4 图1-59所示体系内部是何学变体系，它有\_|

一个多余约束。（广西大学

2010)

图1-59

1-5图1-60所示体系的计算自由度W=

（哈尔滨工业大学2011）

图1-60

1-6图1-61所示体系几何不变且有多余联系。（）（中南大学2012）

图1-61

\subsection*{第 3 页转写}
1-7对图1-62所示体系进行几何组成分析。（哈尔滨工业大学2009）

图1-62

几何不变，三个多余约束。

1-8对图1-63所示体系作几何组成分析。（武汉科技大学2009）

几何不变，没多余约束

图1-63

1-9对图1-64所示杆件体系进行几何构造分析。（武汉理工大学2008）

几何瞬变，一个多余

图1-64

o3是水平方向无穷远点，所以所有水平线都过o3，

那么o102这条水平线也过o3，

因此o1,02,03在一条直线上。

\subsection*{第 4 页转写}
1-10计算图1-65所示体系的计算自由度，试分析其体系的几何组成。（华南理工大

学2012)

12x2- 25: -

图1-65

几何不变，一个多余

1-11对图1-66所示平面体系进行几何组成分析，并指出结构超静定次数。（青岛理

工大学2012)

任何东西都可以作为一个体系

超静定次数是几何不变体系中多余约束数量。

几何不变，2个多余。

图1-66

1- 12

分析图1－67所示体系的几何稳定性，写出分析过程。（东南大学2008)

图1-67

几何不变，1个多余。

\subsection*{第 5 页转写}
1-13

试对图1－68所示体系作几何组成分析。（山东建筑大学2011）

图1-68

几何不变，无多余。

1-14

试对图1-69所示体系作几何组成分析。（山东建筑大学2008）

图1-69

几何不变，无多余。

1- 15

计算图1-70的计算自由度，并作几何构造分析。（中国地震局工程力学研究所

2010)

7x2 - 13 =|

图1-70

常变，无多余约束

\subsection*{第 6 页转写}
1-16对图1-71所示体系进行几何构造分析。（西南交通大学2008）

图1-71

几何瞬变，一个多余

o3是水平方向无穷远点，所以所有水平线都过o3，

那么01o2这条水平线也过o3，

因此o1,02，03在一条直线上。

1-17计算图1-72所示体系的计算自由度，并进行几何构成分析。（北京工业大学

2010)

4-2 -2(-) - 2-2\textasciitilde{} 3= \textasciitilde{}

几何不变，1个多余。

两种分析方法结果不一样，一种不变，一种可变，那么最终结论是不变那种。之所以分析出可变，

是因为没合理利用约束。从我见过的题目和大家常见错误中可以总结出来，一般都是实际上是几

何不变的，有同学却判断出来了几何瞬变，那肯定是瞬变那种错了，针对这种错误，如何自已检

查出来呢，这里有一个风险提示，就是如果你判断出来的瞬变体系有2个多余约束，那你就得特

别警惕了，最好再换换分析方法，看看能不能找到几何不变的分析方法，能找到的话，就是不变

的，找不到的话，那就是瞬变2个多余。

